\subsection*{先删响应的副本}

沿一条独立本地历史，初始余额为 10000，E1 有效。S1 完成 P1/O1/D1，保存完整 $R(S1)$，当次余额为 7000。Alice 收到结果，将 S1、原输入和完整结果写入自己的持久任务记录，再向服务发送确认。这个顺序让设备重启后能够继续认出 P1。

如果 Alice 只在内存里看过成功，就提前确认，随后进程退出，服务可能已经回收结果，客户端也失去自己的一份。确认的含义应是“这项结果已经持久取得，并接受它以后停止回放”。设备持久化失败时，客户端保留原请求，继续按原身份恢复。

确认消息丢失可以重发。服务重复应用确认，保持已确认事实即可。服务收到确认后、回复确认以前退出，客户端再发送同一确认，也只会延续同一回收决定。把确认事实与结果退休一起写入原受理域，副本切换和快照恢复就能保留它。

原结果可能同时在 M 的热缓存、业务表和归档中出现。删除缓存副本以后，如果 C1 的 \code{payment\_result} 或持久结果档案能够重建相同返回，接口可以继续回放。只有删除了最后一份可重建载荷，服务才真正放弃这项能力。物理上删了一行和逻辑上失去哪项保证，应沿实际读取路径判断。

\CTwelveFigure{ch12-original04}{客户端先持久取得结果，再确认回放结束，服务回收载荷并保留身份依据}{fig:ch12-durable-ack}

外部支付沿既有数据归属处理。C1 保存原 \code{provider\_result} 与 \code{payment\_result}，P1 保存业务进度。收款与账务追踪可能要求长期保存 C1 的原证据，因此客户端确认以后，可以先删除另一份为接口加速而保存的大响应。若这些字段本来共用同一存储，相同字节只计算一次节省。后续归档保留原平台范围、\code{pay-C1}、X1 和核查所需关系。

\subsection*{从完整绑定缩到拒绝所需的依据}

第十一章用 $M(S)=(V,X,P,R(S))$ 表示完整身份的绑定：V 决定怎样解释原输入，X 是按该版本规范化后的输入，P 是实际受理的购买，$R(S)$ 是完整原结果。结果退休以后，主方案把它缩为 $M_{\min}(S)$，保留 S、V、X、真实 P、退休状态，以及适用时的原结果档案定位。\code{min} 表示相对于本章完整回放记录的缩减，具体字段由继续提供的能力决定。

保留 X，是为了在同一 S 上继续判断输入是否冲突。X 包含原来按 V 解释后的全部购买输入。只保存一个固定长度摘要，可以减少输入占用，但摘要相等与输入必然相等之间依赖碰撞假设。本章沿用精确输入比较，先把这份成本计入。若接口只承诺拒绝所有退休身份，可以采用更小的已见身份集合，响应同时结束对具体异参的区分。

只保留 S 的已见标记，可以阻止第二次首次执行；加上 P，业务查询就能找到原购买；加上 V/X，服务可以继续报告异参冲突；加上 $R(S)$ 或可重建它的档案，原结果回放才恢复完整。各项能力需要的信息逐层增加。

\CTwelveFigure{ch12-evidence-layers}{从已见标记到完整结果，增加的字段分别支持定位、输入比较与精确回放}{fig:ch12-evidence-layers}

在有效 E1 内，一项主动确认的本地结果可以采用下面的处理。\code{result\_retired} 表示本身份的原结果回放已经结束，调用方应使用自己保存的结果或经过授权的业务查询：

\begin{BookCode}
在同一受理域串行处理 S、原输入和结果确认
若 E1 缺失或准入已关闭：
    返回 operation_expired，不创建购买
否则若 M 或 M_min 已有 S：
    用记录的 V 解释输入，精确比较原 X
    若不同：返回 conflict，保留原绑定
    若有可回放的原 R(S)：返回原 R(S)
    若结果已经退休：返回 result_retired 与已保存的 P
否则：
    进入首次购买仲裁
    共同保存业务变化、完整绑定和原结果

处理结果确认时：
    验证确认者对 S 的权限与确认语义
    仅对已有终态标记 replay_retired
    以 M_min 接替完整绑定，回收可删除的响应载荷
    尚未终结的 C/T 保持原恢复责任
共同发布迁移及其应用边界，提交后才响应
\end{BookCode}

\code{result\_retired} 可以使用单独公布的业务码，例如 409 配合该业务码。它和会话过期的 410 分开：前者说明这一份原结果已经被确认后退休，后者说明旧会话入口资格结束。结果退休若没有 P 可供恢复，响应便省略购买编号。在活动会话中，来件先经过原 V/X 比较，原 S 改换输入时从冲突分支退出。

\subsection*{删除完成映射以后谁接手查找}

M 可以独立建表，也可以作为 P 的受理字段实现。假设 P 已经按完整 S 保存唯一关联、原输入和终态，服务删除独立的 M 行，同时把身份恢复改为查询 P，并用 P 的唯一约束仲裁新受理。这是把信息移到业务数据中。用于恢复的 S 索引及其字段继续占用空间，订单归档以后也要维持可查询性。

若只剩 O1，情况会不同。余额不足的购买没有订单，却同样形成稳定拒绝。只在 \code{orders} 上查找，可能让旧拒绝因订单不存在而重新进入首次判断。外部支付的 A 也可能已经保存 C1/T1、尚未有 O1。按照“有订单就算处理过”的规则清理，会漏掉这些已受理身份。

可以把永久业务记录设计成操作证据，但要覆盖每一种已受理结局，并把已有输入、编号关系及后续首次仲裁接到它上面。对于未完成 C1，平台身份与支付输入先保留到终态证据建立。成功后的 EV1、消费者 I/B 和 F1，又各有消息重复、业务事实仲裁与长期重放的用途，它们按各自事务域的恢复责任决定保留与清理。

\CTwelveFigure{ch12-original05}{O1 只覆盖成交路径，P/C/T 还覆盖拒绝和未决路径，回收按真实关系进行}{fig:ch12-business-evidence}

任意键接口还有一个空间下界。设有效的 E1 允许 n 个候选键，每个键都可能独立已经受理或尚未受理。如果入口以后要对每个键准确区分这两种情况，就存在 $2^n$ 种可能的已见集合。区别这些集合至少需要 n 比特的信息，实际键、输入和索引会占更多空间。把集合移入订单或外部档案，保存和查询这份信息的责任也随之移动。

\subsection*{布隆过滤器能替代哪一段查找}

一种常见的空间优化是布隆过滤器。它用一个初始全为 0 的位数组和若干个散列位置表示已加入的集合。加入一个键时，把它对应的各个位设为 1；查询时检查同样的位置。只要其中有一位为 0，就知道该键未被完整加入；全部为 1 时，其他键也可能恰好共同设过这些位，所以结论是“可能已经存在”。Bloom 原文用这类编码减少常驻存储，让少数误判转入较慢的后续检查~\cite{bloom1970}。

取一个八位数组，位置编号从 0 到 7，每个键对应两个位置。K1 对应 1、4，K2 对应 4、6。依次加入它们以后，1、4、6 为 1，其余为 0。K3 从未加入，若它恰好对应 1、6，查询也会命中。这叫假阳性：一个集合外的键被报告为可能在集合中。K4 对应 0、5，查询遇到 0，可以判断它没有出现在这份完整插入的集合里。

\CTwelveFigure{ch12-bloom-bits}{K1、K2 共同设置位以后，未加入的 K3 也可能命中；命中转向准确索引核对}{fig:ch12-bloom-bits}

这个例子的准确已见集合只有 K1、K2。若入口看到 K3 命中就直接返回“已处理”，会漏掉 Alice 另一项合法购买。若看到 K4 未命中就绕开唯一约束直接扣款，另一请求可能同时通过同一判断。过滤器提供查询提示，持久身份索引负责区分完整 S，首次受理继续在该索引的唯一仲裁下提交。

可以据此安排两条读取路径。命中时先查询准确索引，找到 S 就按原绑定恢复，找不到则参与首次仲裁。未命中时可以少做一次昂贵的档案读取，直接尝试同一个受保护的首次插入；插入冲突后读取胜出记录。无论从哪条路径进入，唯一约束或状态机仲裁都决定哪份 S 成为首次绑定。

\CTwelveFigure{ch12-bloom-route}{过滤器改变读取路径，两种查询结果最后都汇入准确的身份仲裁}{fig:ch12-bloom-route}

“没有假阴性”也有明确前提：所讨论的已见键都已插入这份过滤器，散列规则保持一致，设过的位没有丢失或被清除。在这样的追加结构里，加入 K1 设为 1 的位会一直保留，之后查询 K1 必然命中。删除 K1 时若简单清掉位置 4，K2 也会失去它需要的一位，随后 K2 被错误报告为未见。这叫假阴性。带计数的变体可以记录共享次数，但又要处理计数、重复加入、删除与恢复的一致性。

服务进程重启或切换副本以后，也要检查过滤器来自哪份身份历史。若准确索引已有 S1，过滤器却恢复自较早快照，S1 可能暂时未命中。把这种未命中直接当成未执行，会重现本章的遗忘反例。过滤器可以从完整索引重建，也可以与新增身份共同维护；构建或补齐期间，入口继续走准确索引与唯一仲裁。最终效果由持久绑定决定，过滤器的遗漏只增加读取或冲突处理的工作。

位数固定时，加入的键越多，置 1 的位置通常越多，假阳性随之增加，更多请求回到准确索引。扩大位数组或分段重建可以降低这部分额外查询。若为了永久节省空间而删除最后一份准确已见集合，就只能接受合法新键被误拒的语义，或另行收紧旧输入的准入范围。本章要求两项合法购买准确分开，保留准确身份依据，并在下一节通过确认和会话关闭缩小它的保留范围。
\FloatBarrier
